% !TeX root = ncaa-manuscript-en.tex
\section{Results}\label{sec:results}

The behavioral track first tests training with constraints under common-CNL
evaluation and training-matched evaluation
(Sections~\ref{sec:results-training} and~\ref{sec:results-semantic}), then
examines continuation training (Section~\ref{sec:results-continuation}).
The six-task and rule-pair results assess use of the supplied rule; reasoning,
wording, value, and image controls examine alternative explanations and
robustness. The specification track separately reports fault detection and
candidate-verdict correspondence (Section~\ref{sec:results-checks}). A separate diagnostic tests composition
consistency (Section~\ref{sec:results-composition}).
Section~\ref{sec:results-summary} brings the two tracks together.

Answer-extraction failures are zero in the fresh-adapter, continuation,
regression, six-task, robustness, reasoning-prompt, and exploratory
image/wording evaluations. The pretrained-model evaluations have nonzero
failure rates, reported below and in Supplementary Section S7.

\subsection{Constraint Training with Identical Evaluation Inputs}\label{sec:results-training}

The fresh-adapter study uses two evaluations: inputs matching each training
condition and a common-CNL evaluation for every adapter
(Section~\ref{sec:setup-comparisons}). This subsection emphasizes the latter,
which fixes evaluation information. Section~\ref{sec:results-semantic} reports
the same adapters with training-matched inputs, keeping the effects of input
information and loss design distinct.

\noindent\textbf{Same initialization and evaluation inputs.}
Table~\ref{tab:results-common-rule} compares the four fresh-adapter conditions;
individual-seed values are given in Supplementary Section S6. In the CE-only
comparison, the conditions differ in whether constraint sentences appear in
training. With the same CNL at evaluation, CoC-only-trained adapters violate
the rule in 17.92\% of pedestrian-entry predictions and 6.51\% and 6.25\% of
position-hold predictions with interval maxima and XY coordinates, respectively.
Training with CNL reduces these rates to 1.11\%, 0.00\%, and 1.56\%.
Preferred-selection score increases from 50.14\%, 78.39\%, and 75.52\% to 98.47\%, 96.09\%,
and 93.49\%. With initialization and evaluation information controlled, these
differences support the benefit of including constraint text during training.

\begin{table}[tbp]
\caption{Fresh adapters with identical CNL at evaluation. Cells show preferred-selection score
(first line) and constraint violation rate (second line), in percent, as mean $\pm$
standard deviation over three seeds. Column labels describe training inputs.}
\label{tab:results-common-rule}
\centering\footnotesize
\setlength{\tabcolsep}{3pt}
\begin{tabularx}{\textwidth}{@{}>{\raggedright\arraybackslash}Xcccc@{}}
\toprule
Task & \shortstack{CoC\\CE} & \shortstack{CoC\\CE + semantic} &
\shortstack{CoC + CNL\\CE} & \shortstack{CoC + CNL\\CE + semantic}\\
\midrule
Pedestrian entry & \shortstack{$50.14\pm1.46$\\$17.92\pm8.91$} &
\shortstack{$50.00\pm0.00$\\$0.00\pm0.00$} &
\shortstack{$98.47\pm1.20$\\$1.11\pm1.58$} &
\shortstack{$97.64\pm2.14$\\$0.83\pm1.44$}\\[5pt]
Position hold: maxima & \shortstack{$78.39\pm1.97$\\$6.51\pm3.52$} &
\shortstack{$75.52\pm0.90$\\$0.26\pm0.45$} &
\shortstack{$96.09\pm1.35$\\$0.00\pm0.00$} &
\shortstack{$96.09\pm2.34$\\$0.00\pm0.00$}\\[5pt]
Position hold: XY & \shortstack{$75.52\pm0.45$\\$6.25\pm2.71$} &
\shortstack{$74.74\pm0.45$\\$0.52\pm0.90$} &
\shortstack{$93.49\pm1.19$\\$1.56\pm0.78$} &
\shortstack{$92.71\pm3.61$\\$2.08\pm2.96$}\\
\botrule
\end{tabularx}
\end{table}

CoC-only-trained preferred-selection score remains close to the rule-free input ceilings of
50\% for pedestrian entry and 75\% for position holding, even after CNL is
provided at evaluation. These ceilings apply when rule text is absent
(Section~\ref{sec:setup-comparisons}); they are not mathematical upper bounds
for the common-CNL inputs. Rule-pair accuracy also separates the CE-only
conditions: CoC-only-trained adapters reach 1.94\%, 56.77\%, and 51.56\%,
whereas CNL-trained adapters reach 96.94\%, 92.71\%, and 86.98\%. In these
checkpoints, training with constraint text supports stronger rule-dependent
selection than inference-time rule provision alone.

\noindent\textbf{VLMs without task-specific training.}
Table~\ref{tab:results-pretrained} reports Qwen3-VL-2B-Instruct and
Qwen3-VL-8B-Instruct with the same CoC and CNL inputs but no task-specific
training. Panel (a) reports selection scores together with answer-extraction
failures. At the prespecified 12-token limit, every pretrained 2B response
reaches the cap and yields no extractable label. Panel (b) reports the
post-hoc 512-token evaluation on the three fresh-adapter tasks; extraction
failures remain frequent. Full results and cap-hit counts are reported in
Supplementary Section S7. The 8B model's primary constraint violation rate
is 68.75\% for pedestrian entry, 25.00--50.00\% for the static tasks and
stopping, and 7.03\% and 33.59\% for the two position-hold representations.
Lane changing has no observed violations but 75.00\% preferred-selection score.

Only XY-coordinate responses hit the 8B 12-token limit (26/128). With 512
tokens, pedestrian-entry preferred-selection score / violation rate is
27.50\% / 66.67\%, and XY position holding is 36.72\% / 34.38\%, with
17.19\% extraction failures. Other 8B task scores are unchanged.

\begin{table}[tbp]
\caption{Pretrained 2B and 8B and CNL-trained 2B with CoC + CNL inputs.
(a) Preferred-selection score / constraint violation rate / answer-extraction
failure rate (\%). All use the prespecified 12-token limit; pretrained models
give one greedy response per input, and trained scores are three-seed means.
(b) Post-hoc 512-token results for pretrained 2B on the fresh-adapter tasks.}
\label{tab:results-pretrained}
\centering\footnotesize
\setlength{\tabcolsep}{3pt}
\begin{tabular*}{\textwidth}{@{\extracolsep\fill}lccc@{}}
\toprule
\multicolumn{4}{@{}l}{\emph{(a) Prespecified 12-token evaluation}}\\[3pt]
Task & 2B, no task training & 8B, no task training & 2B, CNL training\\
\midrule
Maximum speed & 0.00 / 0.00 / 100.00 & 60.42 / 39.58 / 0.00 & 100.00 / 0.00 / 0.00\\
Minimum clearance & 0.00 / 0.00 / 100.00 & 54.17 / 45.83 / 0.00 & 100.00 / 0.00 / 0.00\\
Minimum entry delay & 0.00 / 0.00 / 100.00 & 50.00 / 50.00 / 0.00 & 100.00 / 0.00 / 0.00\\
Pedestrian entry & 0.00 / 0.00 / 100.00 & 27.08 / 68.75 / 0.00 & 98.47 / 1.11 / 0.00\\
Stopping & 0.00 / 0.00 / 100.00 & 50.00 / 25.00 / 0.00 & 100.00 / 0.00 / 0.00\\
Lane change & 0.00 / 0.00 / 100.00 & 75.00 / 0.00 / 0.00 & 97.92 / 0.00 / 0.00\\
Position hold: maxima & 0.00 / 0.00 / 100.00 & 70.31 / 7.03 / 0.00 & 96.09 / 0.00 / 0.00\\
Position hold: XY & 0.00 / 0.00 / 100.00 & 35.94 / 33.59 / 20.31 & 93.49 / 1.56 / 0.00\\
\botrule
\end{tabular*}
\par\medskip
\begin{tabular*}{\textwidth}{@{\extracolsep\fill}lccc@{}}
\toprule
\multicolumn{4}{@{}l}{\emph{(b) Pretrained 2B, post-hoc 512-token evaluation (\%)}}\\[3pt]
Task & Preferred-selection score & Violation rate & Extraction failures\\
\midrule
Pedestrian entry & 10.83 & 11.25 & 65.42\\
Position hold: maxima & 14.84 & 0.00 & 79.69\\
Position hold: XY & 7.03 & 7.03 & 75.78\\
\botrule
\end{tabular*}
\par\smallskip\raggedright
Unextractable answers receive zero preferred-selection credit and count as
neither violations nor admissible choices; the pretrained 2B model's zero
12-token violation rate therefore does not indicate compliance. Admissible
lower-scoring choices also receive zero preferred-selection credit, so the
three rates need not sum to 100\%. Trained pedestrian-entry and position-hold
rows use fresh CE adapters; other trained rows use six-task P2 checkpoints.
This is not a matched-training comparison across model sizes. Panel (b)
changes the output limit and extraction protocol (Supplementary Section S7).
\end{table}

The comparison identifies limitations of inference-only rule provision on
these inputs; it does not isolate a model-size effect. Because the pretrained
2B evaluation has substantial extraction failures even at 512 tokens, these
data also cannot isolate whether CoC-only fine-tuning reduced any prior
rule-following ability.

\subsection{Constraints in the Loss and in the Input}\label{sec:results-semantic}

Table~\ref{tab:results-semantic} evaluates the four fresh-adapter conditions
with inputs matching training. Without rule text, semantic loss reduces
pedestrian-entry violations from 19.58\% to 0.00\%, and position-hold
violations from 10.94\% to 1.56\% with maxima and 0.52\% with XY coordinates.
However, preferred-selection score stays near the rule-free ceilings of 50\% and 75\%
(Section~\ref{sec:setup-comparisons}), and rule-pair accuracy remains 0\%
or approximately 50\%.

Half the position-hold pairs share a reference label, permitting up to 75\%
preferred-selection score without distinguishing their rules. Reduced violations therefore do
not by themselves establish rule-dependent choice. The results are consistent
with more conservative selections that are admissible under either rule,
rather than improved discrimination between rules. Providing CNL only at
evaluation does not close the preferred-selection score gap for the semantic-only-trained
condition (Table~\ref{tab:results-common-rule}).

\begin{table}[tbp]
\caption{Fresh adapters evaluated with their training input type. Cells show
preferred-selection score and constraint violation rate (\%, mean $\pm$ seed standard deviation) on the
first two lines; parentheses give rule-pair accuracy (\%). CE denotes
cross-entropy loss; semantic denotes the additional constraint loss defined
in Eq.~\eqref{eq:setup-semantic-loss}.}
\label{tab:results-semantic}
\centering\footnotesize
\setlength{\tabcolsep}{3pt}
\begin{tabularx}{\textwidth}{@{}>{\raggedright\arraybackslash}Xcccc@{}}
\toprule
Task & \begin{tabular}[t]{@{}c@{}}CoC\\CE\end{tabular} &
\begin{tabular}[t]{@{}c@{}}CoC\\CE + semantic\end{tabular} &
\begin{tabular}[t]{@{}c@{}}CoC + CNL\\CE\end{tabular} &
\begin{tabular}[t]{@{}c@{}}CoC + CNL\\CE + semantic\end{tabular}\\
\midrule
Pedestrian entry & \begin{tabular}[t]{@{}c@{}}$49.86\pm0.24$\\$19.58\pm6.14$\\(0.00)\end{tabular} &
\begin{tabular}[t]{@{}c@{}}$50.00\pm0.00$\\$0.00\pm0.00$\\(0.00)\end{tabular} &
\begin{tabular}[t]{@{}c@{}}$98.47\pm1.20$\\$1.11\pm1.58$\\(96.94)\end{tabular} &
\begin{tabular}[t]{@{}c@{}}$97.64\pm2.14$\\$0.83\pm1.44$\\(95.28)\end{tabular}\\[5pt]
Position hold: maxima & \begin{tabular}[t]{@{}c@{}}$75.00\pm0.00$\\$10.94\pm1.35$\\(50.00)\end{tabular} &
\begin{tabular}[t]{@{}c@{}}$74.48\pm0.90$\\$1.56\pm1.56$\\(49.48)\end{tabular} &
\begin{tabular}[t]{@{}c@{}}$96.09\pm1.35$\\$0.00\pm0.00$\\(92.71)\end{tabular} &
\begin{tabular}[t]{@{}c@{}}$96.09\pm2.34$\\$0.00\pm0.00$\\(92.71)\end{tabular}\\[5pt]
Position hold: XY & \begin{tabular}[t]{@{}c@{}}$72.66\pm1.35$\\$10.94\pm2.71$\\(48.44)\end{tabular} &
\begin{tabular}[t]{@{}c@{}}$74.74\pm0.45$\\$0.52\pm0.90$\\(50.00)\end{tabular} &
\begin{tabular}[t]{@{}c@{}}$93.49\pm1.19$\\$1.56\pm0.78$\\(86.98)\end{tabular} &
\begin{tabular}[t]{@{}c@{}}$92.71\pm3.61$\\$2.08\pm2.96$\\(85.42)\end{tabular}\\
\botrule
\end{tabularx}
\end{table}

With CNL training inputs, adding semantic loss gives no consistent further
benefit: preferred-selection score changes by $-0.83$ to $0.00$ percentage points and violation
rate by $-0.28$ to $+0.52$ points. These three-seed estimates do not establish
equivalence between the losses. In this setting, the loss penalty mainly
reduces violations without supplying the rule-dependent input distinction,
whereas training with rule text supports that distinction.

\subsection{Position-Hold Continuation}\label{sec:results-continuation}

The continuation study gives a similar direction of benefit under identical
evaluation inputs (Table~\ref{tab:results-continuation};
Section~\ref{sec:setup-comparisons}). Before position-hold training, adapters
with constraint-bearing histories already make fewer violations when
candidates are represented by interval maxima. For the
2.0 s rule with interval maxima, P0, P1, and P2 violate the rule in 64, 23,
and 20 of 192 predictions. After continuation, the counts are 22 (11.46\%),
1 (0.52\%), and 0. P2 preferred-selection score rises from 81.77\% to 98.44\%, and P1 reaches
the same preferred-selection score with one violation. Every condition receives the same rule
at evaluation, so the differences cannot be attributed to withholding it
from P0 at that stage. They reflect the compared training histories, including
whether constraint sentences were supplied at both training stages.

\begin{table}[tbp]
\caption{Position-hold continuation: strict violation counts and preferred-selection score
(\%) before $\to$ after training. Each cell summarizes 192 predictions (64 rows
over three seeds). Evaluation inputs are identical across conditions.}
\label{tab:results-continuation}
\centering\footnotesize
\setlength{\tabcolsep}{2pt}
\begin{tabular*}{\textwidth}{@{\extracolsep\fill}lcccccc@{}}
\toprule
& \multicolumn{3}{c}{Violations} & \multicolumn{3}{c}{Preferred-selection score (\%)}\\
Representation, hold & P0 & P1 & P2 & P0 & P1 & P2\\
\midrule
Maxima, 0.5 s & $29\to0$ & $13\to1$ & $7\to1$ & $47.92\to67.19$ & $56.25\to89.58$ & $52.08\to94.79$\\
Maxima, 2.0 s & $64\to22$ & $23\to1$ & $20\to0$ & $63.02\to86.98$ & $85.94\to98.44$ & $81.77\to98.44$\\
XY, 0.5 s & $56\to1$ & $49\to5$ & $53\to3$ & $47.92\to61.98$ & $51.56\to81.25$ & $50.00\to90.63$\\
XY, 2.0 s & $121\to23$ & $126\to16$ & $115\to4$ & $32.81\to88.02$ & $33.85\to91.15$ & $39.58\to97.92$\\
\botrule
\end{tabular*}
\end{table}

With XY coordinates and the 2.0 s rule, P2 preferred-selection score rises from 39.58\% to
97.92\%, but four violations remain. Preferred-selection score and constraint compliance are distinct:
under the 0.5 s rule with interval maxima, P0 makes no violations yet reaches
only 67.19\% preferred-selection score, often selecting admissible candidates with lower progress.

On the new position-hold test set, a 0.01 m tolerance leaves violation counts
unchanged for both the fresh-adapter and continuation studies: every prohibited
candidate exceeds the radius by at least 22 mm. The tolerance does affect the
separate regression diagnostic, as reported in Supplementary Section S11.

The earlier regression diagnostic remains difficult after training. Its
XY/2.0 s stratum still yields 106, 79, and 77 violations out of
192 for P0, P1, and P2. The smallest prohibited overshoot is 1 mm in that
diagnostic, compared with 22 mm in the new primary test. The sets are not
pooled; strict and tolerance-based regression results are reported separately
in Supplementary Section S11. The new test supports improved synthetic
selection, with remaining sensitivity to representation and small boundary
overshoots.

\subsection{Use of Supplied Constraints Across Six Tasks}\label{sec:results-six-tasks}

Table~\ref{tab:results-six-tasks} reports the six-task comparison. P0 lacks rule
information and remains near its design-imposed 50\% ceiling. Both P1 and P2
have higher preferred-selection scores and lower constraint violation rates in every task. Their preferred-selection scores
range from 96.81\% to 100\%, with violations at or below 1.39\%.
Pedestrian-entry preferred-selection score exceeds P0 by 46.94 and 48.19 percentage points for
P1 and P2, while violations decrease by 17.78 and 18.33 points. P1 and P2 tie
on the three static tasks and stopping.

\begin{table}[tbp]
\caption{Six-task preferred-selection score / constraint violation rate (\%), three-seed means. These
checkpoints are distinct from the fresh-adapter study in
Tables~\ref{tab:results-common-rule} and~\ref{tab:results-semantic}.}
\label{tab:results-six-tasks}
\centering\footnotesize
\begin{tabular*}{\textwidth}{@{\extracolsep\fill}lccc@{}}
\toprule
Task & P0 & P1 & P2\\
\midrule
Maximum speed & 50.00 / 24.31 & 100.00 / 0.00 & 100.00 / 0.00\\
Minimum clearance & 50.00 / 24.31 & 100.00 / 0.00 & 100.00 / 0.00\\
Minimum entry delay & 50.00 / 24.31 & 100.00 / 0.00 & 100.00 / 0.00\\
Pedestrian entry & 49.86 / 19.17 & 96.81 / 1.39 & 98.06 / 0.83\\
Stopping & 50.00 / 23.61 & 100.00 / 0.00 & 100.00 / 0.00\\
Lane change & 50.00 / 38.19 & 100.00 / 0.00 & 97.92 / 0.00\\
\botrule
\end{tabular*}
\end{table}

P2 lane-change errors select admissible but lower-progress candidates;
constraint-compliant completion is 100\%. P2 pedestrian-entry constraint-compliant completion is 99.17\%.
Supplying constraints in both training and evaluation inputs raises every task
to near-ceiling preferred-selection score. Individual-seed values are given in Supplementary
Section S6.

\subsection{Selection under Rule Changes}\label{sec:results-association}

Changing the rule while holding scene and candidates fixed tests whether
choices change accordingly. Table~\ref{tab:results-association} reports
rule-pair accuracy: P0 reaches 0\% in all three families, whereas P1 and P2
reach 93.61--100\% and 96.11--100\%. Constraint-conditioned models change their
choices with the rule despite unchanged visual input; P0's result reflects its
missing rule information.

In one recorded pedestrian-entry case, clearance
occurs at 11.0 s and a 1.5 s wait permits entry from 12.5 s; P0 selects entry
at 11.5 s, whereas P1 and P2 select the boundary entry at 12.5 s. This and a
P2 violation case are described in Supplementary Section S9.

\begin{table}[tbp]
\caption{Six-task rule-pair accuracy (\%), three-seed means. Static and maneuver values aggregate
their subtasks.}
\label{tab:results-association}
\centering\footnotesize
\begin{tabular*}{\textwidth}{@{\extracolsep\fill}lccc@{}}
\toprule
Family & P0 & P1 & P2\\
\midrule
Pedestrian entry & 0.00 & 93.61 & 96.11\\
Static & 0.00 & 100.00 & 100.00\\
Maneuver & 0.00 & 100.00 & 97.92\\
\botrule
\end{tabular*}
\end{table}

\subsection{General Reasoning-Prompt Controls}\label{sec:results-reasoning}

With the rule omitted, general reasoning prompts leave P0 preferred-selection score and
rule-pair accuracy unchanged
(Section~\ref{sec:setup-additional}). Across the 480 primary-test rows, P0,
Zero-shot CoT, and Plan-and-Solve all reach 49.93\% preferred-selection score, with violation
rates of 23.06\%, 22.64\%, and 22.71\%, and 0\% rule-pair accuracy in every
family. All 546 unique Zero-shot CoT first-stage outputs and 498 of 546 Plan-and-Solve
outputs are short label strings, so the prompts do not establish multistep
reasoning in these label-trained checkpoints. Because identical P0 inputs carry
conflicting reference labels under different rules, no fixed output can exceed
50\% preferred-selection score. Detailed outputs are given in Supplementary Section S8.

\subsection{Representation, Wording, and Unseen Values}\label{sec:results-robustness}

Table~\ref{tab:results-robustness} reports changed wording and numerical values.
Paraphrase rows retain the primary-test cases and rules, so differences
measure wording sensitivity rather than performance on additional independent
scenes. Both constraint conditions substantially outperform P0 on unseen
values. For pedestrian-entry paraphrases, P2 retains 96.94\% preferred-selection score and
1.25\% violations, compared with 90.56\% and 5.00\% for P1.

\begin{table}[tbp]
\caption{Changed wording and values: preferred-selection score / constraint violation rate (\%),
three-seed means. Paraphrases reuse primary-test cases.}
\label{tab:results-robustness}
\centering\footnotesize
\setlength{\tabcolsep}{3pt}
\begin{tabular*}{\textwidth}{@{\extracolsep\fill}lccc@{}}
\toprule
Evaluation & P0 & P1 & P2\\
\midrule
Pedestrian entry: primary & 49.86 / 19.17 & 96.81 / 1.39 & 98.06 / 0.83\\
\quad Paraphrase & 49.86 / 19.17 & 90.56 / 5.00 & 96.94 / 1.25\\
\quad Unseen wait & 49.86 / 18.47 & 96.67 / 2.22 & 96.67 / 1.53\\[3pt]
Static: primary & 50.00 / 24.31 & 100.00 / 0.00 & 100.00 / 0.00\\
\quad Paraphrase & 50.00 / 24.31 & 93.29 / 0.00 & 80.32 / 7.87\\
\quad Unseen values & 50.00 / 22.92 & 99.54 / 0.00 & 97.22 / 0.23\\[3pt]
Maneuver: primary & 50.00 / 30.90 & 100.00 / 0.00 & 98.96 / 0.00\\
\quad Paraphrase & 50.00 / 30.90 & 100.00 / 0.00 & 97.92 / 2.08\\
\quad Unseen values & 50.00 / 30.90 & 100.00 / 0.00 & 98.96 / 0.00\\
\quad Combined difficulty & 50.00 / 30.21 & 94.10 / 0.69 & 91.32 / 1.04\\
\botrule
\end{tabular*}
\end{table}

Static paraphrases reverse the ordering: P2 preferred-selection score falls from 100\% to
80.32\%, with 7.87\% violations, while P1 reaches 93.29\% with no violations.
Both the favorable pedestrian-entry result and the unfavorable static result
are needed to characterize the CNL interface. On the combined maneuver set,
P1 reaches 94.10\% preferred-selection score and 0.69\% violations, compared with 91.32\% and
1.04\% for P2. P2 rule-pair accuracy falls from 97.92\% to 82.64\%.
The combined test cannot identify whether negation, clause order, equality,
or units individually cause the errors.

A smaller exploratory diagnostic (24 rows per wording and seed) shows the same
fragility: a second paraphrase lowers P1 and P2 preferred-selection score to 76.39\% and 77.78\%
(Supplementary Section S10).

Alternative expressions of the same pedestrian-entry rule also perform well
on the primary test (individual seeds in Supplementary Section S6). An
alternative natural-language guard matches P2 preferred-selection score at
98.06\% (1.11\% violations), and logical text, a model-input control rather
than a solver or translator output, reaches 99.17\% (0.14\%). On unseen waits,
the natural-language guard reaches 98.61\% (0.56\%), whereas logical-text
preferred-selection score falls to 81.25\% despite only 0.83\% violations.
Low violation does not imply selecting the highest-progress admissible action.
High primary-test scores thus hold across several rule representations,
while robustness varies with wording and unseen values.

\subsection{Dependence on Visual Input}\label{sec:results-visual}

Image interventions test whether selections change when visual input is
removed or conflicts with the retained case
(Section~\ref{sec:setup-additional}). Although candidate values are supplied
in text, choices depend on the image (three-seed means; full results in
Supplementary Section S10). P2 preferred-selection score against the original
reference falls from 98.61\% to 76.39\% with a blank image and 79.17\% with the
redrawn image, and its choice changes on 25.00\% and 19.44\% of the paired
cases. P1 falls from 100\% to 81.94\% and 87.50\%, with choice changes on
18.06\% and 12.50\%; P0 stays at 50\%.

P2 violations remain at 1.39\% in every image condition. Preferred-selection score loss
therefore does not imply more prohibited choices; changes to admissible but
lower-progress candidates also contribute. The 24 rows comprise three base case configurations, each paired with two
waiting rules and four label rotations (Supplementary Section S10). With
three training seeds, this is an exploratory input-dependence diagnostic.

\subsection{Specification Checking}\label{sec:results-checks}

Table~\ref{tab:results-faults} reports fault injection into the two
specifications used for learning (protocol in Section~\ref{sec:method-checks}).
All 48 valid cases pass, and the full suite detects all 288 supported faults:
the structural, satisfiability, and logic--rule checks detect 192, and sentence
correspondence detects the remaining 96 CNL faults. No fault made a
specification unsatisfiable, so the satisfiability check flagged none. An error
can trigger several layers, so layer counts are not summed. None of the 7,392
probe queries returned UNKNOWN. All 48 jointly wrong evidence--specification
controls remain undetected. If the rule record and its specification field
change to the same wrong value, evidence correspondence and the logic--rule comparison
still agree and the specification remains satisfiable, so detecting such errors
requires evidence independent of the supplied record. An earlier evaluation on
a separate pedestrian-entry specification shows the same correspondence
between fault types and detecting layers (Supplementary Section S4).

\begin{table}[tbp]
\caption{Fault injection into the specifications used for learning. Each row
has one case per base context (24 per task). Counts are accepted cases for
valid specifications and detected cases otherwise; they are cases, not
independent scenes. Detecting-layer names abbreviate the five checks in
Section~\ref{sec:method-checks}.}
\label{tab:results-faults}
\footnotesize
\begin{tabularx}{\linewidth}{@{}>{\raggedright\arraybackslash}p{0.27\linewidth}>{\raggedright\arraybackslash}X>{\raggedright\arraybackslash}X cc@{}}
\toprule
Case & Injection location & Detecting layer & Entry & Hold\\
\midrule
Valid specification & None & None (accepted) & 24/24 & 24/24\\[3pt]
Changed specification duration & Specification; source unchanged & Logic--rule, evidence & 24/24 & 24/24\\[3pt]
Reversed comparison & Logical formula & Logic--rule & 24/24 & 24/24\\[3pt]
Deleted release condition & Logical formula & Logic--rule & 24/24 & 24/24\\[3pt]
Changed unit & Specification & Structural & 24/24 & 24/24\\[3pt]
Changed sentence duration & CNL & Sentence & 24/24 & 24/24\\[3pt]
Changed sentence negation & CNL & Sentence & 24/24 & 24/24\\[3pt]
Jointly wrong evidence and specification & Source record and specification & Not detected & 0/24 & 0/24\\
\midrule
\multicolumn{3}{@{}l}{Supported faults: structural, satisfiability, logic--rule} & 96/144 & 96/144\\[3pt]
\multicolumn{3}{@{}l}{Supported faults: full suite} & 144/144 & 144/144\\
\botrule
\end{tabularx}
\end{table}

Two further checks use different units from the fault-injection case counts.
SMT verdicts agree with the independent integer admissibility calculation on
all 5,376 position-hold candidate checks across training, validation, and test
(Section~\ref{sec:setup-data}). This measures candidate-level correspondence
with the supplied rule.

For the pedestrian-entry specification, four wait values (0.5, 1.5, 0.8,
and 1.8 s) yield 100 queries: four satisfiability and 96 boundary-candidate
queries. All match their expected verdicts: 76 SAT, 24 UNSAT, and no UNKNOWN.
These are counted separately from the 7,392 concrete queries in the
fault-injection evaluation. To illustrate the inclusive boundary, with
clearance at 11.000 s and a 0.500 s wait, entry at 11.499 s is UNSAT,
whereas 11.500 s and 11.501 s are SAT.

The checks identify supported disagreements among specifications, logical
restrictions, sentences, and supplied rules.

\subsection{Constraint-Composition Diagnostic}\label{sec:results-composition}

Table~\ref{tab:results-composition} reports the document-grounded diagnostic
(Sections~\ref{sec:method-composition} and~\ref{sec:setup-composition}). All 176
contexts and individual source groups are satisfiable. Joint checks identify
58 UNSAT and 118 SAT configurations. All 896 subset queries agree with the
independent feasibility reference, with zero UNKNOWN outcomes; independent
numeric checks validate all 832 SAT witnesses. These counts cover distinct
compatibility questions, rather than 58 conflicting real-world policies.
The 58 UNSAT configurations comprise 22 requirement conflicts (12 three-way
braking and 10 yellow-interval), 6 objectives already infeasible without the
delay policy, 22 candidate requests correctly blocked by a source condition,
and 8 crossings rejected only by the omitted-exception fault.

\begin{table}[tbp]
\caption{Document-grounded conditions combined with authored states, models,
objectives, or candidate requests. Counts are diagnostic configurations, not
independent road scenes. Individual groups are SAT. Pairwise detection counts
UNSAT full configurations having an UNSAT two-group subset. The exception
row includes a deliberately faulty transcription and permitted/blocked
candidate controls.}
\label{tab:results-composition}
\centering\footnotesize
\setlength{\tabcolsep}{3pt}
\begin{tabularx}{\linewidth}{@{}Xrrrr@{}}
\toprule
Composition & Cases & Full SAT & Full UNSAT & Pairwise detected\\
\midrule
Suppression + model + stopping goal & 48 & 30 & 18 & 6\\
Yellow stop/proceed feasibility & 64 & 54 & 10 & 10\\
Autonomous-driving interlock + candidate & 32 & 17 & 15 & 15\\
Crossing exception + candidate & 32 & 17 & 15 & 15\\
\midrule
Total & 176 & 118 & 58 & 46\\
\botrule
\end{tabularx}
\end{table}

The suppression family contains 12 three-way conflicts with every pair SAT.
For example, the authored state $v=10$ m/s, $a=5$ m/s$^2$, and $d=15$ m
combines $t_b\ge1$, $x\ge10t_b+10$, and $x\le15$. Each pair is satisfiable;
together they are inconsistent. Removing the delay admits stopping at 10 m
in this simplified model. Six other configurations have an incompatible
stopping model and objective even without suppression; their minimal groups
correctly exclude the delay policy. This distinction prevents attributing
already-infeasible avoidance to a policy interaction. These are hypothetical
states, not the Tempe collision: NTSB reports that avoidance already exceeded
ADS braking specifications at emergency recognition~\cite{NTSB2019Tempe}.
Yellow-interval conflicts arise when neither the stop nor the proceed branch
is feasible under the nominal model (example in Supplementary Section S13).

All 16 interlock hold controls are SAT, while autonomous-driving requests
conflict in the 15 states with at least one blocker. With all three Rule 129
exception conditions true, the faithful encoding accepts a crossing candidate
and the omitted-exception variant rejects it. Both encodings accept the
non-crossing candidate; all 16 such controls are SAT. The exception result
identifies an overly restrictive transcription, not a contradiction in the
original rule. A single UNSAT verdict for a crossing candidate has the same
form whether the candidate is correctly blocked or a permission was lost;
comparing the faithful and omitted-exception encodings on the same candidate
exposes the lost permission. Thus joint checks, read with explicit source
roles and this encoding comparison, distinguish incompatible requirements,
blocked candidates, and a lost permission.
Supplementary Section S13 reports the earlier 123-case synthetic boundary
diagnostic separately.

\subsection{Summary of Findings}\label{sec:results-summary}

Table~\ref{tab:results-summary} collects the central contrast. Every condition
receives the same CoC + CNL inputs at evaluation; only training with constraint
text yields both few violations and high rule-pair accuracy among the
conditions shown.

\begin{table}[htb]
\caption{Central comparison with identical CoC + CNL evaluation inputs:
preferred-selection score / constraint violation rate / rule-pair accuracy (\%). Fine-tuned rows are
three-seed means of fresh 2B adapters (Tables~\ref{tab:results-common-rule}
and~\ref{tab:results-semantic}; Supplementary Section S6); pretrained rows use the prespecified 12-token
limit (Table~\ref{tab:results-pretrained}).}
\label{tab:results-summary}
\centering\footnotesize
\setlength{\tabcolsep}{3pt}
\begin{tabular*}{\textwidth}{@{\extracolsep\fill}lccc@{}}
\toprule
Training & Pedestrian entry & Hold: maxima & Hold: XY\\
\midrule
\multicolumn{4}{@{}l}{\emph{Rules supplied only at evaluation}}\\
None (pretrained 2B) & 0.00 / 0.00 / 0.00\textsuperscript{a} & 0.00 / 0.00 / 0.00\textsuperscript{a} & 0.00 / 0.00 / 0.00\textsuperscript{a}\\
CoC, CE & 50.14 / 17.92 / 1.94 & 78.39 / 6.51 / 56.77 & 75.52 / 6.25 / 51.56\\
CoC, CE + semantic loss & 50.00 / 0.00 / 0.00 & 75.52 / 0.26 / 51.56 & 74.74 / 0.52 / 50.00\\
None (pretrained 8B) & 27.08 / 68.75 / 0.83 & 70.31 / 7.03 / 51.56 & 35.94 / 33.59 / 23.44\\[2pt]
\multicolumn{4}{@{}l}{\emph{Constraint text supplied during training}}\\
CoC + CNL, CE & \textbf{98.47 / 1.11 / 96.94} & \textbf{96.09 / 0.00 / 92.71} &
\textbf{93.49 / 1.56 / 86.98}\\
\botrule
\end{tabular*}
\par\smallskip\raggedright
\textsuperscript{a}Answer-extraction failure is 100\% at the 12-token limit;
the zero violation rate does not indicate compliance (Table~\ref{tab:results-pretrained}).
8B rule-pair accuracy is taken from the evaluation record (Supplementary Section S7).
\end{table}

\begin{itemize}
\item \textbf{Training with constraints.} Relative to CoC-only training, CNL
training lowers violations from 17.92\% to 1.11\% on pedestrian entry and
raises rule-pair accuracy from 1.94\% to 96.94\%
(Section~\ref{sec:results-training}). Continuation training shows the same
direction: with the 2.0 s rule and interval maxima, P0, P1, and P2 retain 22,
1, and 0 violations of 192 (Section~\ref{sec:results-continuation}).
\item \textbf{Alternatives without training-time rule text.} Semantic loss
removes most violations but leaves preferred-selection score near the rule-free
input ceilings of 50\% and 75\% under training-matched evaluation; the pretrained 8B model violates the rule in up to 68.75\% of
predictions; general reasoning prompts leave P0 at 49.93\%
(Sections~\ref{sec:results-training}, \ref{sec:results-semantic},
and~\ref{sec:results-reasoning}).
\item \textbf{Use of supplied constraints.} P1 and P2 reach
96.81--100\% preferred-selection score on the six tasks, compared with about 50\% for P0,
and change their choices when the rule changes
(Sections~\ref{sec:results-six-tasks} and~\ref{sec:results-association}). P2 is
not uniformly better than P1: on static paraphrases, P2 reaches 80.32\% versus
93.29\% for P1 (Section~\ref{sec:results-robustness}).
\item \textbf{Checkability.} The full check suite detects all 288 supported
faults without rejecting the 48 valid cases, but misses all 48 jointly wrong
evidence--specification controls (Section~\ref{sec:results-checks}).
\item \textbf{Composition consistency.} A separate document-grounded diagnostic finds
58 incompatible configurations despite satisfiable individual groups;
12 require checking three sources together, beyond pairwise checks
(Section~\ref{sec:results-composition}).
\end{itemize}

